\graphicspath{{content/chapters/3_methodology/figures/}}

\chapter{Methodology}%
\label{chp:methodology}


\section{Image Retrieval}
%Retrieved images from ReLaion-5B (128 Shards (3.62GB each)) -> Due to space and time constraints reduced to first 4 shards only resulting in a subset of 65,552,809 images. These shards were chosen in order i.e the first 4 shard form the 128 set. This was done as the images themselves don't have any order in how they were saved hence shard selection is irrelevant.
%
%Assessing the origins of the valid images 18.79\% of all valid images are derived from 5 websites these being cdn.shopify.com: 3,053,782, i.pinimg.com: 2,585,519, i.ebayimg.com: 966,820, images-na.ssl-images-amazon.com: 886,054 and thumbs.dreamstime.com: 678,516 in order of contribution.
%
%An image was considered valid so long as the following criteria were met:
%\begin{itemize}
%	\item HTTP Resolution
%	\begin{itemize}
%		\item A GET request to the URL completes successfully.
%		\item The server returns an HTTP success status (2xx).
%		\item The request completes within the defined timeout (10 seconds).
%		\item A browser-like User-Agent header is used.
%		\item A Referer header based on the host domain is supplied.
%		\item No exception is raised during the request (timeout, connection error, SSL error, etc.).
%	\end{itemize}
%	
%	\item Image Decoding
%	\begin{itemize}
%		\item The response content can be opened by PIL.Image.open().
%		\item The image file is not corrupted or truncated.
%		\item The image can be fully loaded via image.load().
%	\end{itemize}
%	
%	\item RGB Conversion
%	\begin{itemize}
%		\item The image is successfully converted to RGB (image.convert("RGB")).
%		\item Images in other formats (RGBA, grayscale, etc.) are normalised to RGB.
%		\item Images that fail conversion are discarded.
%	\end{itemize}
%
%	\item Batch-Level Validation
%	\begin{itemize} 
%		\item If embedding computation fails for a batch:
%		\item The batch is recursively split (binary splitting).
%		\item Only images that individually pass forward propagation are retained.
%		\item Failed images are isolated and excluded.
%	\end{itemize}
%	
%\end{itemize}
%
%Given that the dataset consists of links to images these images were downloaded producing 43,491,938 valid images as some image links were no longer valid. This outliens a decay rate of 0.33 within the subset alinging with teh research in XYZ. 
%
%Given that the focus of this thesis partially revovles around images depciitng people and the spurious features therin, an efficient image search procedure was required. Given that the CLIP embeddings for the LAION-5B dataset are no longer available they had to be re calcualted for the 4 parquet images subsets. This was facilited through the use of the ImageDownloadAndEmbeddings.ipynb file which as per its name processed the indiviudal .parquet donwloaded the images whose links were still valid and producing CLIP embeddings for said iamges using the openai/clip-vit-large-patch14. The choise to impleemnt this search facility was done as CLIP was previously used to facilite this functionality and is widely accapted for this use case furthermore the specific  CLIP model used was chosen as although the original LAION-5B filtering pipeline relied on CLIP ViT-B/32 embeddings, this work employs the higher-capacity ViT-L/14 variant to derive embeddings for retrieval and bias analysis. ViT-L/14 offers increased representational capacity (768-dimensional embeddings, 14×14 patch resolution) and improved semantic alignment, which enhances neighbourhood purity in similarity search and reduces coarse texture bias. Since the objective of this study is not to replicate the LAION filtering process but to analyse representational bias within the dataset, the use of a stronger backbone is methodologically appropriate.
%
%Images were preporcesssed in line with the openai/clip-vit-large-patch14 model which consisted of as per \cite{find-paper}:
%
%\begin{itemize}
%	\item Resize to model’s expected input resolution (typically 224×224).
%	\item Aspect ratio preserved prior to crop.
%	\item Center crop to square.
%	\item Conversion to tensor.
%	\item Pixel value rescaling.
%	\item Channel-wise normalization using CLIP mean/std.
%	\item Conversion to float32 tensors.
%\end{itemize}
%
%This download and embedding process was conducted using an RTX3060ti with 64GB of RAM across the span of 2 to 3 months.





\subsection{Data Acquisition and Subsetting}

Images were retrieved from ReLAION-5B, distributed as 128 shards (approximately 3.62GB each). Due to storage and computational constraints, processing was restricted to the first four shards, yielding an initial subset of 65,552,809 URL--caption pairs. Shards were selected deterministically in index order (i.e., the first four shards). LAION shards are generated during distributed crawling and do not encode semantic ordering; therefore, shard index does not inherently correspond to content category or structure. However, restricting processing to a subset does not guarantee full representativeness. Consequently, all findings derived in this work apply strictly to the processed subset rather than the entire LAION-5B corpus.

\subsection{Image Source Distribution}

An analysis of valid image origins indicated that 18.79\% of all valid images were derived from five domains:
\begin{itemize}
	\item cdn.shopify.com: 3,053,782
	\item i.pinimg.com: 2,585,519
	\item i.ebayimg.com: 966,820
	\item images-na.ssl-images-amazon.com: 886,054
	\item thumbs.dreamstime.com: 678,516
\end{itemize}

This concentration suggests strong domain centralisation, with large commercial platforms contributing disproportionately. Such concentration may introduce commercial or e-commerce visual bias, as product imagery and stock photography may be overrepresented relative to smaller independent domains.

\subsection{Image Validity Criteria}

Since ReLAION-5B contains URLs rather than stored image binaries, each URL was resolved and downloaded. An image was considered valid only if all of the following conditions were satisfied:

\paragraph{HTTP Resolution}
\begin{itemize}
	\item A GET request completed successfully.
	\item The server returned an HTTP success status (2xx).
	\item The request completed within a 10-second timeout.
	\item A browser-like \texttt{User-Agent} header was used.
	\item A domain-based \texttt{Referer} header was supplied.
	\item No network exception (timeout, DNS failure, SSL error, connection error) occurred.
\end{itemize}

\paragraph{Image Decoding}
\begin{itemize}
	\item The response content could be opened using \texttt{PIL.Image.open()}.
	\item The file was not corrupted or truncated.
	\item The image successfully executed \texttt{image.load()}.
\end{itemize}

\paragraph{RGB Conversion}
\begin{itemize}
	\item The image was successfully converted to RGB using \texttt{image.convert("RGB")}.
	\item Non-RGB formats (e.g., RGBA, grayscale) were normalised to RGB.
	\item Images that failed conversion were discarded.
\end{itemize}

\paragraph{Batch-Level Validation During Embedding}
\begin{itemize}
	\item If embedding computation failed for a batch, recursive binary splitting was applied.
	\item Images that individually passed forward propagation were retained.
	\item Failed samples were isolated and excluded.
\end{itemize}

\subsection{Decay Rate}

From 65,552,809 attempted URLs, 43,491,938 valid images were successfully retrieved. The decay rate was computed as:

\[
\text{Decay Rate} = 1 - \frac{43,491,938}{65,552,809} \approx 0.33
\]

Decay includes HTTP failures, decoding failures, and conversion failures. This level of link attrition aligns with prior findings in large-scale web datasets and reflects natural web degradation over time. However, link decay may introduce survivorship bias if certain domains or content types disappear at disproportionate rates.

\subsection{Image Preprocessing}

Images were preprocessed using the canonical pipeline associated with \texttt{openai/clip-vit-large-patch14} via the HuggingFace \texttt{AutoProcessor}. Preprocessing steps included:

\begin{itemize}
	\item Resize to the model's expected resolution (typically 224$\times$224).
	\item Preservation of aspect ratio prior to cropping.
	\item Center crop to square.
	\item Conversion to tensor representation.
	\item Pixel value rescaling.
	\item Channel-wise normalisation using CLIP mean and standard deviation.
	\item Conversion to \texttt{float32} tensors.
\end{itemize}

No data augmentation was applied. All preprocessing was deterministic to ensure reproducibility.

\subsection{Embedding Computation}

CLIP embeddings were recalculated for all valid images, as original LAION-5B embeddings are no longer publicly available. Embedding extraction was performed using:

\begin{itemize}
	\item Model: \texttt{openai/clip-vit-large-patch14}
	\item Pretraining corpus: OpenAI WebImageText
	\item Embedding dimensionality: 768
	\item Patch resolution: 14$\times$14
\end{itemize}

Images were processed in batches (batch size = 200) and forwarded through the CLIP image encoder using \texttt{model.get\_image\_features(...)} under \texttt{torch.no\_grad()} in evaluation mode. GPU acceleration was utilised when available.

Each embedding vector $\mathbf{e}$ was L2-normalised:

\[
\mathbf{e}_{norm} = \frac{\mathbf{e}}{\|\mathbf{e}\|_2}
\]

This ensures unit-length representations and enables cosine similarity search via dot-product equivalence. Embeddings were stored in \texttt{float32} format and written incrementally to Parquet files with checkpointing for fault tolerance and resumability.

\subsection{Model Selection Justification}

While the original LAION-5B filtering pipeline relied on CLIP ViT-B/32 embeddings, this work employs the higher-capacity ViT-L/14 variant for retrieval and bias analysis. ViT-L/14 provides increased representational capacity (768-dimensional embeddings), finer spatial granularity (14$\times$14 patches versus 32$\times$32), and improved semantic alignment in retrieval benchmarks. As the objective of this thesis is not to replicate LAION filtering thresholds but to analyse representational bias, the use of a higher-capacity backbone is methodologically appropriate. Nonetheless, embeddings inherit inductive biases from their pretraining corpus, and retrieval behaviour remains model-dependent.

\subsection{Computational Environment}

The download and embedding pipeline was executed on an NVIDIA RTX 3060 Ti GPU with 64GB system RAM. Processing spanned approximately two to three months. The extended duration may introduce minor temporal drift in URL availability.

\subsection{Ethical Considerations}

\begin{itemize}
	\item Images were accessed exclusively via publicly available URLs.
	\item No authentication or access restrictions were bypassed.
	\item Images were used strictly for academic research purposes.
	\item No redistribution of image binaries was performed.
	\item Analysis was conducted at aggregate statistical levels.
	\item No attempt was made to identify individuals.
	\item No facial recognition or re-identification techniques were applied.
	\item Where possible, only embeddings were retained for analytical tasks.
	\item The research objective is to analyse bias rather than amplify or deploy it.
\end{itemize}

\subsection{Limitations}

\begin{itemize}
	\item Only 4 of 128 shards were processed; findings are limited to this subset.
	\item Shard selection does not guarantee full representativeness.
	\item Approximately 33\% of URLs were unavailable, introducing survivorship bias.
	\item 18.79\% of valid images originated from five domains, indicating domain concentration.
	\item Embeddings reflect biases inherent to CLIP pretraining data.
	\item Image resizing to 224$\times$224 may reduce fine-grained demographic cues.
	\item A 10-second timeout may classify slow servers as invalid.
	\item No geographic replication of downloads was performed.
	\item Results are encoder-dependent and not model-agnostic.
	\item No additional content filtering was applied prior to embedding extraction.
\end{itemize}

%Following initial image retrieavl and embeddding creation images needed to be selected upon which spurious feature investiaget would be carried out. To facilitate this in line with prior research the descion was made to focus on profession related images, to this extent a set of 200 varied professions was constructed aligned with prior research as outlined in Section \ref{subsec:occupations}. [Should the profession list be outlined here or in the Appendix?] To facilitate image retrieval il line with these professions the previosuly derived embeddings were used, and two approaches were tested a brute-force appraoch where each mebddings was comapred to the text clip embedding retrieving images having the highest simialrty however this appraoch proved too time consuming having O(n) time complexity n = number of valid images. Hence FAISS was adopted, as outlined in Section \ref{subsec:faiss} it offers various degrees by which search can be carried out, in the case of this approach exact search through IndexFlatIP was utilised as although performing exact search time complexity was reduced in comapriosn to the brute force approach. 
%
%This was facilitated through the use of FAISS indices constructed across the valid image set embeddings, making large-scale similarity-based retrieval computationally tractable. In line with this, for each profession the decision was made to retrieve the top 1,000 valid images, resulting in 200,000 images in total. This number was chosen to ensure statistically stable prevalence estimates of spurious features within each profession.
%
%Under a binomial modelling assumption for spurious feature presence (i.e., feature present vs. absent), a sample size of 1,000 per profession yields a worst-case 95\% margin of error of approximately ±3.1\% for proportion estimates. This level of precision ensures that observed differences across professions are not attributable to sampling noise but instead reflect substantive variation in feature prevalence.
%
%Furthermore, selecting an equal number of images per profession produces a balanced evaluation set. This design prevents more frequently occurring professions in the underlying dataset from disproportionately influencing aggregate statistics, thereby enabling controlled cross-profession comparisons. When considered collectively, the 200,000-image balanced subset provides highly precise aggregate estimates (worst-case ±0.22\% at 95\% confidence), although such aggregate measures are interpreted as reflecting an equal-weighted average across professions rather than the natural profession distribution in LAION-5B.
%
%TO facitie image retrieval a prompt tempalte "A photo of a/an \{profession\}" was utilised retrieving the top 125k images with a cosine similarity threshold of 0.15 or above. These were arbitrarily chosen to ensure that invalid or redudndet images where fully removed whilst reducing the number of images requiring further processing. Further more the following keyowrds "Cartoon", "NSFW", "Sex", "Naked", "Object", "Sign", "Logo" "Document", "Paper" and "Page" were used as negative prompts to avoid retrieivng innapropriate images as well as images depicting documents which initialy were compsoing the majority of images retrieve fro professions such as Lawyer. This furhter processing consisted of retrieving the top 10k images per profession from within the 125k subset having the highest similarity score whilst also clearly depciitng indiviudals. TO faciltiet tis latter checks 2 models were applied these being Yolo and Yolo-face the check consisitied of processing each image across the two models ensuring that a face is clearly visible which was to be used later adn the person detection served as an extra failsafe to ensure that faces where not incorreclty detected in images where no people were present. (expalin this in more detail aligned with the code) THe two models used were yolov8n-face.pt & yolov8n.pt respectively. These were utilsied as they produced the best performance when comapred across other models variants such as MediaPipe Face Detection, SCRFD and RetinaFace when tested on the WIDER Face Validation Image dataset as seen in Table \ref{add_table}. Yolov8 for person detection was not evalued due to tis widepsread sue and accepted viability and performance.
%
%Following CLIP retrieval, images were filtered using YOLO-based person and face detection combined with FaceMesh segmentation. Across all professions, a total of 25,125,679 candidate images were processed. Of these, 8,420,684 images satisfied all acceptance criteria, yielding an overall acceptance rate of 33.51%.
%
%Acceptance rates varied across professions, with a mean of 33.51\%, median of 33.67\%, and standard deviation of 10.71\%. The lowest acceptance rate observed was 9.55\%, while the highest was 59.21\%. This variability reflects differences in how strongly particular profession prompts correspond to images containing clearly detectable single individuals.
%
%The overall rejection rate of approximately 66\% indicates that CLIP similarity retrieval alone frequently returns images lacking clear human structure, reinforcing the necessity of explicit person and face validation prior to demographic analysis
%
%In addition to the above a random subset of 10k images which were not used for the professiion 200k set were randomly chosen from the LAION image set to serve as a test set when applying the techniques outlined in \cite{} to ensure that the spuriosu features identified don't solely apply the profession related images howver are present across the entire dataset.

\subsection{Profession-Conditioned Subset Construction}

Following large-scale image retrieval and embedding computation, a structured subset of images was required upon which spurious feature analysis would be conducted. In line with prior work examining occupational representation and stereotype formation in web-scale datasets (Section~\ref{subsec:occupations}), the decision was made to focus on profession-conditioned imagery. A set of 200 diverse professions was constructed and aligned with established ISCO-based occupational categories. For clarity and readability, the complete profession list is provided in the Appendix.

Previously computed CLIP ViT-L/14 image embeddings (768-dimensional, L2-normalised) were used to retrieve profession-aligned images. Two retrieval strategies were evaluated. First, a brute-force similarity computation was implemented, wherein each image embedding was compared against the corresponding CLIP text embedding using cosine similarity. This approach has time complexity $\mathcal{O}(n)$, where $n$ denotes the number of valid images (over 43 million). Although exact, this method proved computationally infeasible at scale.

Consequently, FAISS was adopted (Section~\ref{subsec:faiss}). An \texttt{IndexFlatIP} index was constructed over the full embedding set. Since embeddings were L2-normalised, inner product similarity is equivalent to cosine similarity. While \texttt{IndexFlatIP} performs exact search, FAISS provides highly optimised vectorised implementations, resulting in substantial practical speed improvements relative to naïve brute-force computation.

\subsection{Retrieval Parameters and Prompt Design}

Profession retrieval was performed using the template:

\[
\texttt{"A photo of a/an \{profession\}"}
\]

This phrasing was selected to encourage naturalistic depictions of individuals rather than symbolic or abstract representations.

For each profession:
\begin{itemize}
	\item The top 125{,}000 images were initially retrieved.
	\item A cosine similarity threshold of 0.15 was applied.
	\item From this candidate pool, the top 1{,}000 highest-scoring structurally valid images were selected.
\end{itemize}

The similarity threshold was empirically chosen to remove near-random matches while maintaining a sufficiently large candidate pool to enable robust filtering.

Additionally, the following exclusion keywords were applied during retrieval filtering: \emph{Cartoon, NSFW, Sex, Naked, Object, Sign, Logo, Document, Paper, Page}. These exclusions were introduced after exploratory analysis revealed frequent retrieval of document-based or non-photographic imagery (e.g., legal contracts for ``lawyer''), which are unsuitable for demographic analysis.

\subsection{Sample Size Justification}

For each of the 200 professions, 1{,}000 images were retained, yielding a balanced subset of 200{,}000 images.

Under a binomial assumption for spurious feature presence (feature present vs.\ absent), a sample size of 1{,}000 per profession yields a worst-case 95\% margin of error of approximately $\pm 3.1\%$ for proportion estimates (maximal at $p = 0.5$). This ensures that observed differences in feature prevalence across professions are unlikely to arise from sampling noise.

Selecting an equal number of images per profession produces a balanced evaluation set. This design prevents more frequently occurring professions in the underlying dataset from disproportionately influencing aggregate statistics, thereby enabling controlled cross-profession comparison. When aggregated across all professions, the 200{,}000-image subset yields highly precise overall prevalence estimates (worst-case $\pm 0.22\%$ at 95\% confidence), though such aggregate measures reflect an equal-weighted average rather than the natural profession distribution within LAION-5B.

\subsection{Human-Centric Structural Filtering}

CLIP retrieval does not guarantee that returned images contain clearly identifiable individuals. Therefore, a multi-stage structural validation pipeline was applied using YOLO-based detection models and MediaPipe FaceMesh segmentation.

Filtering was performed using:
\begin{itemize}
	\item \texttt{yolov8n-face.pt} (face detection),
	\item \texttt{yolov8n.pt} (person detection).
\end{itemize}

Face detection required:
\begin{itemize}
	\item Exactly one detected face,
	\item Confidence $\geq 0.6$.
\end{itemize}

Person detection required:
\begin{itemize}
	\item At least one detected person (class ID = 0),
	\item Confidence $\geq 0.5$.
\end{itemize}

To ensure that detected faces corresponded to detected persons, spatial consistency constraints were enforced:
\begin{itemize}
	\item Face centroid must lie within a person bounding box \textbf{or}
	\item Intersection-over-Union (IoU) $\geq 0.3$.
\end{itemize}

Additionally, the face-to-person area ratio was required to exceed 0.02 to prevent degenerate detections. Following successful validation, MediaPipe FaceMesh segmentation was applied to extract a refined facial region. Images failing any stage were marked invalid.

Model selection was informed by empirical evaluation on the WIDER Face validation dataset (see Table~\ref{add_table}), where YOLOv8-face demonstrated favourable performance relative to MediaPipe Face Detection, SCRFD, and RetinaFace. Person detection was included as an additional structural safeguard to prevent false-positive face detections in non-human contexts.

\subsection{Filtering Outcomes and Acceptance Rate Distribution}

Across all professions, a total of 25{,}125{,}679 candidate images were processed following initial retrieval. Of these, 8{,}420{,}684 images satisfied all structural acceptance criteria, yielding an overall acceptance rate of 33.51\%.

Acceptance rates varied substantially across professions. The mean acceptance rate was 33.51\%, the median 33.67\%, and the standard deviation 10.71\%. The lowest observed acceptance rate was 9.55\%, while the highest was 59.21\%.

Figure~\ref{fig:acceptance_histogram} presents the distribution of acceptance rates across professions. The histogram demonstrates that structural validity following CLIP retrieval is highly profession-dependent. Certain professions systematically yield retrieval results that more frequently depict clear single-person imagery, while others produce noisier or object-dominated results. This variability indicates that structural filtering does not affect all professions uniformly and may introduce differential pruning effects across categories.

The overall rejection rate of approximately 66\% further illustrates that similarity-based retrieval alone is insufficient for ensuring demographic suitability. Explicit structural validation is therefore necessary prior to fairness or spurious feature analysis.

\subsection{Duplicate Handling}

Images with identical embeddings were removed to prevent duplication within the profession-conditioned subset. Cross-profession duplication was theoretically possible if an image ranked highly for multiple profession prompts; however, such instances were empirically rare and were retained when present, as they represent semantically relevant examples for each respective query.

\subsection{Random Control Subset}

To evaluate whether identified spurious features were specific to profession-conditioned retrieval or reflective of broader dataset characteristics, an additional control subset of 10{,}000 randomly sampled images was constructed. These images were drawn without replacement from the processed ReLAION subset, excluding images used in the profession-conditioned dataset.

Sampling was seed-controlled to ensure reproducibility. This control group serves to distinguish retrieval-induced bias from underlying dataset bias and to assess whether observed spurious correlations generalise beyond occupation-aligned subsets.

In line with the goals of this reaserch paper generated images derived from models trained on this dataset were required. Several models were considered and tested these being SDv1.5, SDv1.5 + Speed Lora, SD v2.1, SDXL, SDXL+Speed Lora, SDXL-Turbo, SSD-1B,  SSD-1B+Speed Lora, PixArt-Alpha, PixArt-Alpha LCM, Segmind Vega + Vega RT Lora. These models were compared to one anotehr in terms of inference speed, image quality was not considered in relation to identifying the model to utilise as this is not relevant to the question tackled in this thesis, serving only as a secondary factor for model selection. In terms of inference time as per the table \ref{tab:generative-model-inference-time} the SDv1.5 + Speed Lora model was chosen as it was amongst the fastest models whilst producing coherent images.



\subsection{Image Generation Pipeline}

In order to analyse representational bias at the generative stage, synthetic images were produced using diffusion-based text-to-image models trained on LAION-scale datasets. Multiple models were evaluated, including SDv1.5, SDv1.5 + Speed LoRA, SDv1.2, SDXL, SDXL + Speed LoRA, SDXL-Turbo, SSD-1B, SSD-1B + Speed LoRA, PixArt-Alpha, PixArt-Alpha LCM, and Segmind Vega (with Vega RT LoRA). 

Models were compared primarily on inference speed rather than perceptual image quality. Since the objective of this thesis is to analyse structural and demographic bias rather than aesthetic fidelity, image quality was considered a secondary factor. As shown in Table~\ref{tab:generative-model-inference-time}, SDv1.5 with Latent Consistency Model (LCM) LoRA was selected. This configuration achieved one of the lowest inference times while maintaining structurally coherent outputs suitable for downstream validation.

Stable Diffusion v1.5 was loaded with LCM LoRA weights and fused for fast sampling using a four-step inference schedule. The scheduler was replaced with an LCM scheduler to enable low-step generation while maintaining semantic consistency. Generation was performed in half-precision (FP16) on GPU.

\subsection{Prompt Design and Justification}

The final prompt used for profession-conditioned generation was:

\begin{quote}
	\texttt{"full-body realistic photo of a \{profession\}, standing, centered, in a professional setting appropriate for their occupation, sharp focus, 35mm lens, natural professional lighting"}
\end{quote}

Each component was selected for functional and fairness-related reasons.

\paragraph{Full-body framing.}
Diffusion models trained on LAION-scale data exhibit a strong bias toward portrait-style close-ups. Explicitly requesting ``full-body'' reduces framing bias and ensures compatibility with structural validation constraints requiring a full detected person bounding box.

\paragraph{Standing, centered.}
Explicit pose and spatial positioning tokens reduce multi-subject ambiguity and mitigate edge-cropping artefacts. This increases the probability of successful face–person spatial association during validation.

\paragraph{Professional setting appropriate for their occupation.}
Rather than specifying explicit environments (e.g., hospital, courtroom), which are known to amplify demographic priors, an abstract contextual phrase was used. This preserves semantic grounding while limiting overfitting to stereotypical scenes.

\paragraph{Sharp focus.}
Low-step LCM sampling can introduce blur. High-frequency detail cues such as ``sharp focus'' improve edge clarity and increase downstream detector reliability.

\paragraph{35mm lens.}
Camera tokens act as implicit geometric priors influencing subject distance and body proportions. A 35mm focal length promotes natural perspective and full-body framing without excessive distortion.

\paragraph{Natural professional lighting.}
Lighting significantly influences aesthetic regression scores and object detection reliability. Neutral professional lighting reduces noise artefacts and improves face and person detection confidence.

\subsection{Generation with Structural Validation}

Images were generated for each of the 200 professions using the above prompt template. For each profession, 1,000 valid images were required.

The generation pipeline incorporated asynchronous structural validation (see \texttt{ImageGeneration.py}~\cite{turn2file0}):

\begin{itemize}
	\item YOLO person detection (confidence $\geq 0.4$),
	\item YOLO face detection (confidence $\geq 0.4$, exactly one face required),
	\item Spatial face–person consistency via centroid inclusion or bounding box expansion,
	\item Face-to-person area ratio thresholding,
	\item MediaPipe FaceMesh segmentation for anatomical refinement,
	\item CLIP-based aesthetic scoring (minimum score threshold = 3.0).
\end{itemize}

An image was accepted only if all criteria were satisfied. Images failing any stage were stored in an invalid directory with the rejection reason encoded in the filename.
% Valid images were saved alongside segmented facial crops and associated metadata, including bounding boxes and aesthetic scores.

%The validation process was executed asynchronously using a multi-threaded worker pool to prevent CPU-bound detection tasks from stalling GPU-based image generation. This pipelined design ensured efficient utilisation of computational resources while maintaining strict structural constraints.

\subsection{Generation Without Validation (Control)}

To assess generative behaviour absent structural filtering, an additional control set of 10,000 images was generated using blank prompts and no validation constraints (see \texttt{ImageGeneration\_NoValidation.py}~\cite{turn2file1}). 

These images were sampled using identical inference parameters but without YOLO-based filtering, spatial constraints, or aesthetic thresholding. This control group enables comparison between structurally constrained and unconstrained generative outputs, isolating the effect of validation logic on observed representational patterns.

\subsection{Rationale for Structural Validation in Generation}

Unlike retrieval from ReLAION-5B, generative models can produce:
\begin{itemize}
	\item Multiple individuals,
	\item Truncated bodies,
	\item Disconnected faces,
	\item Anatomically inconsistent compositions,
	\item Low-quality or blurred renderings.
\end{itemize}

Since downstream demographic and spurious feature analysis requires clearly identifiable single individuals, structural validation is necessary to ensure comparability between retrieved and generated subsets. 

This design ensures that differences observed between real and synthetic images reflect model behaviour rather than artefacts of composition, cropping, or detection instability.

\subsection{Structural Acceptance Rates for Generated Images}

Following generation and validation, acceptance rates were computed per profession in an identical manner to the retrieval analysis. 

Across 200 professions, the mean structural acceptance rate for generated images was 66.38\%, with a median of 71.00\% and standard deviation of 16.60\%. Acceptance rates ranged from approximately 10\% to over 90\%.

Figure~\ref{fig:sd_acceptance_histogram} illustrates the distribution of acceptance rates across professions. Compared to retrieval from ReLAION-5B (mean acceptance 33.51\%), generated images exhibit substantially higher structural validity. This indicates that diffusion outputs, when guided by explicit full-body and framing constraints, are more likely to satisfy single-person detection and spatial consistency requirements than real-world web imagery.

However, variability across professions remains pronounced. Certain occupations consistently produce high acceptance rates, suggesting that the diffusion model readily associates these roles with canonical single-person depictions. In contrast, other professions exhibit lower structural acceptance, potentially reflecting greater compositional ambiguity or weaker learned visual priors within the model.

The higher mean acceptance rate suggests that structural constraints imposed via prompt engineering and diffusion priors systematically bias generation toward detector-compatible compositions. This difference between retrieved and generated subsets must be considered when comparing downstream demographic or spurious feature prevalence.


\section{}

%Methodology Note — CFG Scale Justification
%What to mention:
%State that CFG was fixed at s = 1 across all image generation experiments and explain what this means technically — that at s = 1 no guidance amplification is applied, so the model produces purely conditional outputs without any exaggeration of the prompt signal.
%Why to mention it:
%Reviewers familiar with SD v1.5 will notice s = 1 is atypical (the standard default is 7–7.5) and may flag it as a limitation if left unexplained. Getting ahead of this turns a potential weakness into a strength.
%How to frame it:
%Present it as a deliberate and principled methodological choice rather than a default. The reasoning is that fixing s = 1 isolates bias attributable to the model's learned training distribution from bias amplified at inference time. Since your thesis investigates training-data-derived spurious features, this is the most appropriate setting.
%The key argument to make:
%Any demographic bias identified under s = 1 represents a conservative lower bound. At the guidance scales used in real-world deployment (s = 7–7.5), the same biases would be amplified further. This strengthens rather than weakens your findings — if bias is detectable at its minimum, it will be more pronounced in practice, reinforcing the urgency of addressing it at the level of model weights and training data rather than inference settings.
%One line summary of the argument:
%Finding bias at s = 1 means the bias is structural, not incidental.